\section{Engenharia de Dados para Inteligência Artificial}

\textbf{OBJETIVO GERAL:}

Desenvolver competências para coletar, integrar, transformar, armazenar e disponibilizar dados estruturados e não estruturados, aplicando princípios de Engenharia de Dados para suportar aplicações de Inteligência Artificial, Machine Learning e Analytics.

\textbf{CARGA HORÁRIA:} 120 horas

\textbf{FUNÇÃO:}

Desenvolver sistemas computacionais baseados em Inteligência Artificial, Análise Preditiva e Ecossistemas IoT, atendendo às normas de qualidade corporativas, robustez metodológica, governança de dados e arquitetura segura.

\textbf{SUBFUNÇÕES:}

\begin{itemize}
\item Coletar dados provenientes de diferentes fontes.
\item Preparar bases de dados para aplicações de Inteligência Artificial.
\item Desenvolver pipelines de processamento de dados.
\item Armazenar dados em bancos relacionais, NoSQL e Data Lakes.
\item Garantir qualidade, governança e segurança dos dados.
\end{itemize}

\textbf{PADRÕES DE DESEMPENHO:}

\begin{itemize}
\item Identificar diferentes fontes de dados.
\item Desenvolver processos de ETL e ELT.
\item Integrar dados estruturados e não estruturados.
\item Validar qualidade dos dados.
\item Modelar bancos de dados para aplicações analíticas.
\item Automatizar pipelines de dados.
\item Aplicar princípios de governança.
\item Garantir segurança durante o processamento dos dados.
\end{itemize}

\textbf{CAPACIDADES TÉCNICAS:}

\textbf{Coleta de Dados}

\begin{itemize}
\item Identificar diferentes fontes de dados.
\item Coletar dados utilizando APIs.
\item Integrar bases de dados heterogêneas.
\item Automatizar processos de ingestão.
\end{itemize}

\textbf{Preparação de Dados}

\begin{itemize}
\item Limpar bases de dados.
\item Tratar valores ausentes.
\item Padronizar dados.
\item Validar consistência das informações.
\end{itemize}

\textbf{Modelagem de Dados}

\begin{itemize}
\item Modelar bancos relacionais.
\item Modelar bancos NoSQL.
\item Organizar Data Lakes.
\item Estruturar Data Warehouses.
\end{itemize}

\textbf{Pipelines de Dados}

\begin{itemize}
\item Desenvolver processos ETL.
\item Desenvolver processos ELT.
\item Automatizar fluxos de processamento.
\item Monitorar pipelines.
\end{itemize}

\textbf{Governança}

\begin{itemize}
\item Aplicar princípios de qualidade de dados.
\item Garantir rastreabilidade.
\item Aplicar políticas de segurança.
\item Documentar pipelines.
\end{itemize}

\textbf{CONHECIMENTOS:}

\textbf{Fundamentos de Engenharia de Dados}

\begin{itemize}
\item Engenharia de Dados
\item Ciclo de Vida dos Dados
\item Fontes de Dados
\item Tipos de Dados
\item Qualidade dos Dados
\end{itemize}

\textbf{Coleta e Ingestão}

\begin{itemize}
\item APIs REST
\item Arquivos CSV
\item JSON
\item XML
\item Web Scraping (conceitos)
\end{itemize}

\textbf{Preparação de Dados}

\begin{itemize}
\item Limpeza de Dados
\item Tratamento de Valores Nulos
\item Normalização
\item Padronização
\item Transformação de Dados
\end{itemize}

\textbf{Bancos de Dados Relacionais}

\begin{itemize}
\item PostgreSQL
\item Modelagem Relacional
\item SQL
\item Índices
\item Consultas
\end{itemize}

\textbf{Bancos NoSQL}

\begin{itemize}
\item MongoDB
\item Documentos
\item Coleções
\item Modelagem NoSQL
\end{itemize}

\textbf{Data Warehouse e Data Lake}

\begin{itemize}
\item Data Warehouse
\item Data Lake
\item Lakehouse
\item Arquiteturas Analíticas
\end{itemize}

\textbf{Pipelines}

\begin{itemize}
\item ETL
\item ELT
\item Orquestração
\item Automatização
\item Monitoramento
\end{itemize}

\textbf{Governança de Dados}

\begin{itemize}
\item Qualidade dos Dados
\item Catálogo de Dados
\item Linhagem de Dados
\item Metadados
\item Governança
\end{itemize}

\textbf{Segurança}

\begin{itemize}
\item LGPD
\item Controle de Acesso
\item Criptografia
\item Anonimização
\item Auditoria
\end{itemize}

\textbf{CAPACIDADES SOCIOEMOCIONAIS:}

\begin{itemize}
\item Demonstrar organização no tratamento dos dados.
\item Zelar pela qualidade das informações processadas.
\item Avaliar criticamente inconsistências nos dados.
\item Trabalhar colaborativamente na construção de pipelines.
\item Produzir documentação técnica dos processos.
\item Compartilhar boas práticas de engenharia de dados.
\item Demonstrar responsabilidade na manipulação de dados sensíveis.
\item Atuar observando princípios éticos e legais relacionados aos dados.
\end{itemize}

\textbf{AMBIENTES PEDAGÓGICOS:}

\begin{itemize}
\item Laboratório de Informática.
\item Laboratório de Banco de Dados.
\item Laboratório de Desenvolvimento de Software.
\item Ambiente Virtual de Aprendizagem.
\item Ambiente Cloud para processamento de dados.
\end{itemize}

\textbf{EQUIPAMENTOS E RECURSOS:}

\textbf{Hardware}

\begin{itemize}
\item Computadores com acesso à Internet.
\item Servidor para banco de dados.
\item Projetor multimídia.
\end{itemize}

\textbf{Software}

\begin{itemize}
\item Python
\item PostgreSQL
\item MongoDB
\item DBeaver
\item VS Code
\item Git
\item Docker
\item Jupyter Notebook
\item Google Colab
\end{itemize}
